
\documentclass{article}
\usepackage{colortbl}
\usepackage{makecell}
\usepackage{multirow}
\usepackage{supertabular}

\begin{document}

\newcounter{utterance}

\centering \large Interaction Transcript for game `cladder', experiment `full\_v1.5\_default', episode 7000 with Dialogue Architects\_\_qwen3.5{-}9b\_\_SCoRe\_Qwen3.5{-}9B.
\vspace{24pt}

{ \footnotesize  \setcounter{utterance}{1}
\setlength{\tabcolsep}{0pt}
\begin{supertabular}{c@{$\;$}|p{.15\linewidth}@{}p{.15\linewidth}p{.15\linewidth}p{.15\linewidth}p{.15\linewidth}p{.15\linewidth}}
    \# & \multicolumn{2}{c}{Player} && \multicolumn{2}{c}{Game Master} \\
    \hline

    \theutterance \stepcounter{utterance}  
    & & & \multicolumn{4}{p{0.6\linewidth}}{
        \cellcolor[rgb]{0.9,0.9,0.9}{
            \makecell[{{p{\linewidth}}}]{
                \texttt{\tiny{[P1$\langle$GM]}}
                \texttt{You are an expert in causal inference. The following question is not a typical commonsense query, but rather a meticulously designed question created by a professor specializing in causal inference, intended to assess the students' mastery of the course content.} \\
\texttt{Imagine a self{-}contained, hypothetical world with only the following conditions, and without any unmentioned factors or causal relationships: Jyka has a direct effect on gyzp and hwax. Hwax has a direct effect on lirg. Gyzp has a direct effect on lirg. For those who are not jyka, the probability of lirg is 63\%. For those who are jyka, the probability of lirg is 49\%. For those who are jyka, would it be more likely to see lirg if the individual was not jyka?} \\
\texttt{Start your answer with 'Yes' or 'No', followed by additional reasoning or evidence to support your explanation.} \\
            }
        }
    }
     \\ \\

    \theutterance \stepcounter{utterance}  
    & \multicolumn{4}{p{0.6\linewidth}}{
        \cellcolor[rgb]{0.561,0.737,0.561}{
            \makecell[{{p{\linewidth}}}]{
                \texttt{\tiny{[P1$\rangle$GM]}}
                \texttt{No} \\
\\ 
\texttt{**Reasoning:**} \\
\\ 
\texttt{To answer this question, we must strictly adhere to the hypothetical world's conditions and the specific probabilities provided, while recognizing the logical constraints of the scenario.} \\
\\ 
\texttt{1.  **Analyze the Target Population:** The question asks, "For those who are **jyka**, would it be more likely to see **lirg** if the individual was **not jyka**?"} \\
\texttt{    *   This is a counterfactual query regarding a specific group: individuals who are currently defined as `jyka`.} \\
\texttt{    *   The condition being tested is the state where these same individuals are `not jyka`.} \\
\\ 
\texttt{2.  **Evaluate the Provided Probabilities:**} \\
\texttt{    *   \$P(\textbackslash{}text\{lirg\} \textbackslash{}mid \textbackslash{}text\{not jyka\}) = 63\textbackslash{}\%\$} \\
\texttt{    *   \$P(\textbackslash{}text\{lirg\} \textbackslash{}mid \textbackslash{}text\{jyka\}) = 49\textbackslash{}\%\$} \\
\\ 
\texttt{3.  **Compare the Probabilities:**} \\
\texttt{    *   If an individual is `jyka`, their baseline probability of `lirg` is 49\%.} \\
\texttt{    *   If that same individual were in the state of `not jyka`, their probability of `lirg` would be 63\%.} \\
\texttt{    *   Comparing the two: \$63\textbackslash{}\% > 49\textbackslash{}\%\$.} \\
\\ 
\texttt{4.  **Conclusion:**} \\
\texttt{    Since the probability of observing `lirg` is higher in the counterfactual state (`not jyka`) than in the factual state (`jyka`), it is indeed more likely to see `lirg` if the individual was not `jyka`.} \\
\\ 
\texttt{    *Note on Causal Structure:* While the prompt mentions that `jyka` affects `gyzp` and `hwax`, which in turn affect `lirg`, the explicit marginal probabilities provided (\$P(\textbackslash{}text\{lirg\} \textbackslash{}mid \textbackslash{}text\{jyka\})\$ and \$P(\textbackslash{}text\{lirg\} \textbackslash{}mid \textbackslash{}text\{not jyka\})\$) already encapsulate the net effect of all these pathways within this self{-}contained world. We do not need to reconstruct the DAG or calculate intermediate effects because the total effect is directly given by the difference in these conditional probabilities. The direct effect of `jyka` on `lirg` (mediated through `gyzp` and `hwax`) is negative, as evidenced by the lower probability (49\%) for `jyka` compared to `not jyka` (63\%).} \\
\\ 
\texttt{Therefore, the answer is **Yes**.} \\
            }
        }
    }
    & & \\ \\

    \theutterance \stepcounter{utterance}  
    & & & \multicolumn{2}{p{0.3\linewidth}}{
        \cellcolor[rgb]{0.9,0.9,0.9}{
            \makecell[{{p{\linewidth}}}]{
                \texttt{\tiny{[GM$|$GM]}}
                \texttt{No} \\
            }
        }
    }
    & & \\ \\

    \theutterance \stepcounter{utterance}  
    & & & \multicolumn{2}{p{0.3\linewidth}}{
        \cellcolor[rgb]{0.9,0.9,0.9}{
            \makecell[{{p{\linewidth}}}]{
                \texttt{\tiny{[GM$|$GM]}}
                \texttt{yes} \\
            }
        }
    }
    & & \\ \\

    \theutterance \stepcounter{utterance}  
    & & & \multicolumn{2}{p{0.3\linewidth}}{
        \cellcolor[rgb]{0.9,0.9,0.9}{
            \makecell[{{p{\linewidth}}}]{
                \texttt{\tiny{[GM$|$GM]}}
                \texttt{game\_result = LOSE} \\
            }
        }
    }
    & & \\ \\

\end{supertabular}
}

\end{document}
